\section{Rank-three markings}\label{rank-three-markings}

Throughout, \(x^g=g^{-1}xg\) and \([x,y]=x^{-1}y^{-1}xy\).

A rank-three involutory marking is an ordered generating tuple of nonidentity involutions. For a subset
\(I\) of its indices, \(G_I\) denotes the subgroup on those markers. The
intersection property is \(G_I\cap G_J=G_{I\cap J}\). We retain the standard
regular-polytope correspondence. The following standard rank-three criterion is included with its elementary proof.

\textbf{Rank-three test.} For three distinct involutions with commuting
endpoints and adjacent products of order at least three, IP is equivalent
to \(G_{01}\cap G_{12}=G_1\).

\emph{Proof.} Necessity is immediate. For sufficiency, the endpoint subgroup
is \(G_{02}=\{1,r_0,r_2,r_0r_2\}\). If \(G_{02}\cap G_{01}\) contains
\(r_2\) or \(r_0r_2\), then \(r_2\in G_{01}\cap G_{12}=G_1\), contradicting
distinctness. Thus this intersection is \(G_0\); dually
\(G_{02}\cap G_{12}=G_2\). These equalities exclude every generator from
the subgroup on the other two. Intersections involving singletons now
follow, and intersections involving the empty or full index set are
immediate. This covers all subset pairs. \(\square\)
